\section{Conclusion}

We present \textsc{CrackID}, which consists of 140 photographs from 17 walls, along with a protocol for retrieval, same-defect decisions, and open-set identification with rejection.

The main finding is negative.
For RQ1, ten zero-shot crop matchers exceed chance Rank-1 in point estimates, but with a validation-calibrated threshold their $F_1$ stays near the constant baseline.
For RQ2, wall context appears to contribute more than crack appearance, and the geometric reference improves open-set identification only on the 48\% of alignable valid pairs, with results that may depend on annotation provenance.
For RQ3, the small synthetic revisit probe provides no evidence of multi-visit robustness; OSNet+ctx leads on retrieval and open-set identification, whereas \hybrid{} is marginally better on assignment $F_1$.

The study focuses on single-visit photographs: the figures are point estimates derived from 19 matchable test identities, grouped by wall and evaluated against labels from a single annotator.
The next step is a second acquisition session to build a cross-visit benchmark, followed by ticket-level evaluation.
